\subsection*{把这组课程当成一条重新学习的路线}
这条目录位置是课程公告，不是一节包含新定理的数学课。它在教材中的作用，是说明后续基础主题怎样连接：二次方程教我们用中心化简结构，单位圆教我们用运动解释三角，复数把旋转加入乘法，对数再把乘法翻译成加法。

可以先读本组，再进入线性代数与微积分；也可以遇到复指数或对数困惑时回来补基础。学习顺序不应只按公式难度排，而应看后一概念需要借用前一概念的哪些解释。例如复数幂需要先知道指数的周期性与对数分支，否则容易只剩形式计算。

\subsection*{怎样把观看变成能独立复述的理解}
每次暂停后，用自己的例子解释正在解决的问题。不要先问“这一页有哪些公式”，而要问“为什么需要这个新表示，上一种表示哪里不够方便”。如果无法说出引入符号的原因，通常说明还没有把推导连接起来。

做练习时先保留预测，再看答案。错误预测可以告诉你混淆了哪些对象：把坐标当向量，把有限步近似当极限，或把必要条件当充分条件。修正这类误解比多抄一遍结论更有价值。

本书把连续讲解放在前面，把定义回顾、算例、易错点和 Lean 边界放在后面。先沿着叙述重建思路，再用回顾定位概念，最后核对证明中的假设。公告页本身没有机器定理，这种区分也是认真记录学习范围的一部分。
